% ================================================================
% 《分析讲义》(LECTURE NOTES IN ANALYSIS, 2011, Sergiu Klainerman)
% 中文译本主模板 —— 基于数译(mathtranslations.org)排版模板修改
% 使用 XeLaTeX 编译两次以上(生成目录、交叉引用、术语索引页码).
% 依赖: logo.pdf
%
% 相对原模板的修改:
%  1. 定理类环境(定义/引理/定理/推论/命题/注/例/问题)共享同一计数器,
%     编号显示为"节号.序号"(与原书一致;无节之章与章引言为"0.x").
%  2. 新增 \mtsubsec{标题}: 节内带编号子标题,推进共享计数器(如"3.5. 能量色散").
%  3. 新增 \mtsubsecn{数字}{标题}: 独立编号的子标题,不推进计数器(第 1 章用).
%  4. 新增无编号环境 definition* / lemma* / theorem* / corollary* /
%     proposition* / remark* / example* / problem*(原书未编号对象用).
% ================================================================

\documentclass[UTF8,11pt,fontset=none]{ctexart}

% ---------- 页面、字体与数学 ----------
\usepackage[paper=a4paper,left=2.5cm,right=2.5cm,top=2.8cm,bottom=2.8cm]{geometry}
\usepackage{fontspec}
\usepackage{xeCJK}
\usepackage{amsmath,amsthm,amssymb,mathrsfs,mathtools}
\usepackage{graphicx,tikz,tikz-cd,titlesec}
\usepackage{array,booktabs,longtable}
\usepackage{etoolbox}
\usepackage[dvipsnames]{xcolor}
\usepackage[
    colorlinks=true,
    linkcolor=MidnightBlue,
    citecolor=BrickRed,
    urlcolor=MidnightBlue,
    linktoc=all,
    bookmarksnumbered=true
]{hyperref}

% 英文: Computer Modern Unicode Serif
\setmainfont{cmunrm.otf}[
  BoldFont=cmunbx.otf,
  ItalicFont=cmunti.otf,
  BoldItalicFont=cmunbi.otf
]

% 中文: 正文宋体;专用术语楷体;定理类环境仿宋
\setCJKmainfont{FandolSong-Regular.otf}[BoldFont=FandolSong-Bold.otf]
\newCJKfontfamily\songfont{FandolSong-Regular.otf}[BoldFont=FandolSong-Bold.otf]
\newCJKfontfamily\termfont{FandolKai-Regular.otf}
\newCJKfontfamily\theoremfont[AutoFakeBold=2]{FandolFang-Regular.otf}
\setCJKsansfont{FandolHei-Regular.otf}[BoldFont=FandolHei-Bold.otf]
\setCJKmonofont{FandolFang-Regular.otf}

\pagestyle{plain}
\setlength{\parindent}{2em}
\setlength{\parskip}{0pt}
\linespread{1.16}

% ---------- 译本元信息 ----------
\newcommand{\BookTitleCN}{分析讲义}
\newcommand{\BookTitleEN}{LECTURE NOTES IN ANALYSIS (2011)}
\newcommand{\OriginalAuthor}{Sergiu Klainerman}
\newcommand{\OriginalEdition}{2011 年讲义(普林斯顿大学)}
\newcommand{\OriginalPublisher}{Princeton University}
\newcommand{\OriginalYear}{2011}
\newcommand{\Translator}{数译}
\newcommand{\ModelUsed}{DeepSeek-V4-Flash}
\newcommand{\TranslationDate}{18 Aug 2026}

% ================================================================
%                         术语系统
% ================================================================
\newcommand{\termcn}[1]{{\termfont #1}}

\makeatletter
\newcommand{\MTTermList}{}
\newcommand{\MTRegisterTerm}[3]{%
  \g@addto@macro\MTTermList{\MTTermRow{#1}{#2}{#3}}%
}
\newcommand{\MTTermRow}[3]{%
  {\termfont #1} & #2 & \hyperref[#3]{\pageref*{#3}} \\
}

\newcommand{\newterm}[3]{%
  \@ifundefined{MT@seen@#1}{%
    \expandafter\gdef\csname MT@seen@#1\endcsname{}%
    \phantomsection
    \label{term:#1}%
    \protected@write\@auxout{}{%
      \string\MTRegisterTerm{\unexpanded{#2}}{\unexpanded{#3}}{term:#1}%
    }%
  }{}%
  {\termfont #2}（#3）%
}
\makeatother

\newcommand{\printterminology}{%
  \clearpage
  \phantomsection
  \section*{术语索引}
  \addcontentsline{toc}{section}{术语索引}

  \vspace{0.5cm}
  \begingroup
  \renewcommand{\arraystretch}{1.28}
  \begin{longtable}{>{\raggedright\arraybackslash}p{0.28\textwidth}>{\raggedright\arraybackslash}p{0.49\textwidth}>{\centering\arraybackslash}p{0.12\textwidth}}
    \toprule
    \textbf{中文术语} & \textbf{英文术语} & \textbf{页码} \\
    \midrule
    \endfirsthead
    \toprule
    \textbf{中文术语} & \textbf{英文术语} & \textbf{页码} \\
    \midrule
    \endhead
    \MTTermList
    \bottomrule
  \end{longtable}
  \endgroup
}

% ================================================================
%                       定理类环境
% ================================================================
% 定义、引理、定理、推论、命题、注: 仿宋正文
% 例、问题: 宋体正文
% 环境标题后使用中文冒号.

\newtheoremstyle{mtfang}%
  {6pt}{6pt}%
  {\theoremfont}%
  {0pt}%
  {\songfont\bfseries}%
  {.}%
  {0.5em}%
  {}%

\newtheoremstyle{mtsong}%
  {6pt}{6pt}%
  {\songfont}%
  {0pt}%
  {\songfont\bfseries}%
  {.}%
  {0.5em}%
  {}%

\theoremstyle{mtfang}
% 所有定理类环境共享同一个计数器 theorem,按 subsection(原书"节")重置.
\newtheorem{theorem}{定理}[subsection]
\newtheorem{definition}[theorem]{定义}
\newtheorem{lemma}[theorem]{引理}
\newtheorem{corollary}[theorem]{推论}
\newtheorem{proposition}[theorem]{命题}
\newtheorem{remark}[theorem]{注}

\theoremstyle{mtsong}
\newtheorem{example}[theorem]{例}
\newtheorem{problem}[theorem]{问题}

% 无编号版本(原书未编号的定义、定理、注等)
\newtheorem*{theorem*}{定理}
\newtheorem*{definition*}{定义}
\newtheorem*{lemma*}{引理}
\newtheorem*{corollary*}{推论}
\newtheorem*{proposition*}{命题}
\newtheorem*{remark*}{注}
\newtheorem*{example*}{例}
\newtheorem*{problem*}{问题}

% 编号显示与原书一致: "节号.序号"(如 3.4);无节章/章引言为 "0.x"
\renewcommand{\thetheorem}{\arabic{subsection}.\arabic{theorem}}
% 节标题显示为 "1." "2." 形式(与原书一致,不带章号)
\renewcommand{\thesubsection}{\arabic{subsection}.}

% 节内带编号子标题: 推进共享计数器并显示 "节号.序号 标题"
\newcommand{\mtsubsec}[1]{%
  \par\medskip\noindent\refstepcounter{theorem}%
  {\songfont\bfseries\arabic{subsection}.\arabic{theorem}\quad #1}\par\smallskip
}
% 节内独立编号子标题: 只显示 "数字. 标题",不推进计数器(第 1 章情形)
% 用法: \mtsubsecn{数字}{标题}(数字为 0 或空时只显示标题)
\newcommand{\mtsubsecn}[2]{%
  \par\medskip\noindent%
  {\songfont\bfseries\ifx\relax#1\relax\else #1.\quad\fi#2}\par\smallskip
}

% ---------- 按类型区分的引用命令 ----------
% 所有定理类环境共享同一计数器,故 \autoref 无法区分类型(一律显示"定理").
% 引用时用下列命令: \thmref{...} → "定理 N", \lemref{...} → "引理 N" 等.
\newcommand{\thmref}[1]{定理~\ref{#1}}
\newcommand{\lemref}[1]{引理~\ref{#1}}
\newcommand{\defref}[1]{定义~\ref{#1}}
\newcommand{\propref}[1]{命题~\ref{#1}}
\newcommand{\corref}[1]{推论~\ref{#1}}
\newcommand{\remref}[1]{注~\ref{#1}}
\newcommand{\exref}[1]{例~\ref{#1}}
\newcommand{\probref}[1]{问题~\ref{#1}}

% autoref 名称
\providecommand{\definitionautorefname}{定义}
\renewcommand{\definitionautorefname}{定义}
\providecommand{\lemmaautorefname}{引理}
\renewcommand{\lemmaautorefname}{引理}
\providecommand{\theoremautorefname}{定理}
\renewcommand{\theoremautorefname}{定理}
\providecommand{\corollaryautorefname}{推论}
\renewcommand{\corollaryautorefname}{推论}
\providecommand{\propositionautorefname}{命题}
\renewcommand{\propositionautorefname}{命题}
\providecommand{\exampleautorefname}{例}
\renewcommand{\exampleautorefname}{例}
\providecommand{\remarkautorefname}{注}
\renewcommand{\remarkautorefname}{注}
\providecommand{\problemautorefname}{问题}
\renewcommand{\problemautorefname}{问题}
\providecommand{\exerciseautorefname}{习题}
\renewcommand{\exerciseautorefname}{习题}

% 只有"例"和"注"以菱形表示环境结束.
\newcommand{\envenddiamond}{%
  \ifhmode
    \unskip\nobreak\hfill\ensuremath{\diamond}%
  \else
    \noindent\hfill\ensuremath{\diamond}%
  \fi
  \par
}
\AtEndEnvironment{example}{\envenddiamond}
\AtEndEnvironment{remark}{\envenddiamond}

% ---------- proof 样式 ----------
\makeatletter
\renewcommand{\proofname}{证明}
\renewenvironment{proof}[1][\proofname]{\par
  \pushQED{\qed}%
  \normalfont\songfont
  \topsep6\p@\@plus6\p@\relax
  \trivlist
  \item[\hskip\labelsep\normalfont\songfont\bfseries #1：]\ignorespaces
}{%
  \popQED\endtrivlist\@endpefalse
}
\makeatother
\renewcommand{\qedsymbol}{\ensuremath{\square}}

% ================================================================
%                     长证明双向链接
% ================================================================
\newcommand{\longprooflink}[2]{%
  \par\smallskip
  \hypertarget{longproof-source:#1}{}%
  \begin{flushright}
    \fbox{\hyperref[longproof:#1]{\textcolor{MidnightBlue}{#2（第~\pageref*{longproof:#1}~页）}}}%
  \end{flushright}
  \smallskip
}

\newenvironment{longproof}[2]{%
  \def\CurrentLongProofKey{#1}%
  \phantomsection
  \hypertarget{longproof:#1}{}%
  \label{longproof:#1}%
  \begin{proof}[证明（#2）]
}{%
  \end{proof}
  \begin{flushright}
    \fbox{\hyperlink{longproof-source:\CurrentLongProofKey}{\textcolor{MidnightBlue}{返回原处}}}%
  \end{flushright}
  \medskip
}

% ================================================================
%                   Chapter 风格的 section
% ================================================================
\newlength{\chapterNumSpace}
\setlength{\chapterNumSpace}{1.4cm}
\newif\ifInAppendix
\InAppendixfalse
\newif\ifInAppendixToc
\InAppendixTocfalse

\newlength{\tocSectionLabelWidth}
\setlength{\tocSectionLabelWidth}{4.2em}
\newlength{\tocAppendixGap}
\setlength{\tocAppendixGap}{0.2em}
\newlength{\tocSubIndent}
\setlength{\tocSubIndent}{1.2em}
\newlength{\tocSubNumWidth}
\setlength{\tocSubNumWidth}{3.0em}

\newcommand{\TocSectionNumber}[1]{%
  \ifInAppendixToc
    附录\hspace{\tocAppendixGap}#1%
  \else
    第~#1~章%
  \fi
}

\makeatletter
\renewcommand*\l@section[2]{%
  \ifnum\c@tocdepth>\z@
    \addvspace{0.55em}%
    \begingroup
      \parindent\z@
      \rightskip\@pnumwidth
      \parfillskip-\@pnumwidth
      \leavevmode
      \bfseries
      \setlength{\@tempdima}{\tocSectionLabelWidth}%
      \advance\leftskip\@tempdima
      \hskip-\leftskip
      \def\numberline##1{\makebox[\@tempdima][l]{\TocSectionNumber{##1}}}%
      #1\nobreak\hfill\nobreak\makebox[\@pnumwidth][r]{#2}\par
    \endgroup
  \fi
}
\renewcommand*\l@subsection{\@dottedtocline{2}{\tocSubIndent}{\tocSubNumWidth}}
\makeatother

\let\OriginalAppendix\appendix
\renewcommand{\appendix}{%
  \OriginalAppendix
  \InAppendixtrue
  \addtocontents{toc}{\protect\InAppendixToctrue}%
}

\newcommand{\ChapterSectionLabel}{%
  \thispagestyle{plain}%
  \vspace*{-0.5cm}%
  \makebox[\textwidth][r]{%
    {\fontsize{18pt}{18pt}\selectfont\scshape
      \ifInAppendix Appendix\else Chapter\fi
    }\enspace
    {\fontsize{62pt}{62pt}\selectfont\bfseries\thesection}%
  }%
  \par\vspace{-0.25cm}%
  \noindent\rule{\dimexpr\textwidth-\chapterNumSpace\relax}{0.45pt}%
}

\titleformat{\section}[display]
  {\normalfont}
  {\ChapterSectionLabel}
  {0cm}
  {\raggedleft\rightskip=\chapterNumSpace\fontsize{30pt}{35pt}\selectfont\bfseries}
  [\vspace{0.3cm}]

\newcommand{\sectionbreak}{\clearpage}

\titleformat{name=\section,numberless}[block]
  {\normalfont\Large\bfseries\centering}
  {}
  {0pt}
  {}

\titlespacing{\section}{0pt}{0pt}{0pt}
\ctexset{subsection={format=\large\bfseries}}

% 独立大标题(用于"第 1 部分"等),进入目录
\newcommand{\bookpart}[1]{%
  \clearpage
  \phantomsection
  \addcontentsline{toc}{section}{#1}%
  \begin{center}
    {\Large\bfseries #1}
  \end{center}
  \vspace{1.2cm}
}

% ================================================================
%                     习题与答案双向链接
% ================================================================
\newcounter{exercise}[subsection]
\renewcommand{\theexercise}{\arabic{subsection}.\arabic{exercise}}
\providecommand{\theHexercise}{\theexercise}
\newlength{\exerciseLabelWidth}
\newlength{\answerLabelWidth}
\newcounter{answeritem}

\makeatletter
\newcommand{\maybehyperref}[2]{\@ifundefined{r@#1}{#2}{\hyperref[#1]{#2}}}

\newcommand{\exercise@item}{%
  \par\smallskip
  \refstepcounter{exercise}%
  \edef\exercise@number{\theexercise}%
  \edef\exercise@label{exe:\exercise@number}%
  \edef\exercise@answerlabel{ans:\exercise@number}%
  \expandafter\label\expandafter{\exercise@label}%
  \settowidth{\exerciseLabelWidth}{\textbf{\theexercise.}\enspace}%
  \hangindent=\exerciseLabelWidth
  \hangafter=1
  \noindent
  \maybehyperref{\exercise@answerlabel}{\textbf{\theexercise.}}\enspace
  \ignorespaces
}

\newenvironment{exercises}{%
  \par\medskip
  \noindent{\large\bfseries 习题}\par\smallskip
  \begingroup
  \setlength{\parindent}{0pt}
  \renewcommand{\item}{\exercise@item}%
}{%
  \par\endgroup\medskip
}

\newcommand{\answer@item}{%
  \par\smallskip
  \stepcounter{answeritem}%
  \edef\answer@number{\answer@prefix.\arabic{answeritem}}%
  \edef\answer@label{ans:\answer@number}%
  \edef\answer@exerciselabel{exe:\answer@number}%
  \phantomsection
  \protected@edef\@currentlabel{\answer@number}%
  \expandafter\label\expandafter{\answer@label}%
  \settowidth{\answerLabelWidth}{\textbf{\answer@number.}\enspace}%
  \hangindent=\answerLabelWidth
  \hangafter=1
  \noindent
  \maybehyperref{\answer@exerciselabel}{\textbf{\answer@number.}}\enspace
  \ignorespaces
}

\newenvironment{answers}[1]{%
  \par\medskip
  \def\answer@prefix{#1}%
  \setcounter{answeritem}{0}%
  \begingroup
  \setlength{\parindent}{0pt}
  \renewcommand{\item}{\answer@item}%
}{%
  \par\endgroup\medskip
}
\makeatother

% ---------- 参考文献环境 ----------
\newenvironment{mybibliography}[1]{%
  \begingroup
  \renewcommand{\section}[2]{}%
  \begin{thebibliography}{#1}%
}{%
  \end{thebibliography}%
  \endgroup
}
